\section{A geometric probe detects every literal variational source defect}
\label{sec:r12-single-probe-normal-residual}

Retain the strict original analytic box and background-extension
hypotheses \textup{(ED1)}.  The background $S$ is arbitrary:
it is the original symmetric middle-positive separately closed
$(p,p)$-current, with $q=n-p\ge1$.  Set $d=q+1$ and keep
the unchanged strong classical-approximation cone $\mathscr C$.
Let $\psi_0$ be the single geometry-dependent bounded actual
probe of \textup{(UP2)--(UP5)}.  Its construction precedes
the choice of $S$, and it is a finite-mass, strictly negative,
continuous zero-boundary actual potential for every allowed $S$.

The result below concerns compact original smooth directions.
It does not claim that $\mathscr C$ has nonempty Banach
interior, or that a pointwise lower barrier permits subtraction
of that barrier inside the psh cone.

\begin{lemma}[A simultaneous two-sided original test probe]
\label{lem:r12-two-sided-geometric-probe}
For every real original compact smooth scalar $\phi$ on $V$
there is $\epsilon_\phi>0$, depending only on the geometric
probe and the chosen actual test, such that
\[
             \psi_0+t\phi\in\mathscr T_0(S,V)
                  \qquad(|t|\le\epsilon_\phi)
 \tag{NP1}
\]
for every allowed background $S$ and every remaining degree.
The actual functions in \textup{(NP1)} are bounded continuous
classical psh, strictly negative, of ordinary zero boundary
value and finite total top mass.  They agree with $\psi_0$
outside one compact set.  Every actual moment refinement of
$\phi$ is permitted.

On one fixed finite original moment presentation, compact
support and smooth partition, one may choose a constant
$c_{\mathfrak a}>0$ such that
\[
       \epsilon_\phi\ge
           \frac{c_{\mathfrak a}}{\max(1,\|\phi\|_{2,\mathfrak a})}.
 \tag{NP2}
\]
Here $\|\cdot\|_{2,\mathfrak a}$ is the sum of the
usual real $C^2$ norms of the smooth representatives in
that fixed presentation.  No presentation-independent
global norm on all original tests is asserted.
\end{lemma}

\noindent\textbf{Status: \texttt{NEW\_PROVED}} for the
simultaneous probe.  The finite smooth perturbation argument
is the frozen explicit classical compensator construction in
R11 \textup{(AE5)} and the proof of \textup{(HC15)}, rather
than an inference from positivity against one current.

\begin{proof}
Take a finite original smooth partition of the support and
write $\phi=\sum_{l\in F}\phi_l$, with each $\phi_l$
supported in the interior of a compact moment core used
in \textup{(UP1)}.  Include the finitely many moment
functions of the partition itself in those presentations.
The countable exact moment replacement and its closure under
finite refinements allow these cores and compensators to
be selected from the family defining $\psi_0$.
If an index repeats, first combine its finitely many
localized test functions.  The construction then gives
\[
 \psi_0=\sum_{l\in F}c_l\psi_l+r_F,
 \qquad \psi_l=a_l-2\text{ near }\operatorname{supp}\phi_l,
 \tag{NP3}
\]
where $r_F$ is the uniformly convergent sum of the remaining
nonpositive actual psh terms.  It is itself actual bounded
classical psh and belongs to every printed finite power class;
the zero residual is harmless.

On the relevant smaller original moment stage, the explicit
convex-flat numerator and denominator summands in $a_l$
have a strictly positive Hessian in the finite selected
logarithmic directions.  The real Hessian coefficients of
$\phi_l$ are bounded on their fixed compact support.
Consequently there is $\delta_l>0$ such that
\[
                     a_l+\tau\phi_l
             \text{ is actual classical psh near that support}
                  \qquad(|\tau|\le\delta_l).
 \tag{NP4}
\]
This is exactly the frozen explicit smooth perturbation
argument: the selected convex-flat terms dominate both
signs of the Hessian of the smooth real representative,
and the remaining compensator terms are already classical
psh.  Convex composition is used in those actual finite
unit charts, followed by ordinary-open gluing.  No claim
that a current inequality for one background implies
classical psh is made.

Since $\phi_l$ vanishes off its compact support,
$\psi_l+\tau\phi_l$ is therefore classical on all of
$V$ for these small $\tau$.  Decrease $\delta_l$ if
necessary so that it is strictly negative: the continuous
$-\psi_l$ has a positive minimum on that support.
Choose $\epsilon_\phi>0$ with
$\epsilon_\phi/c_l\le\delta_l$ for every $l\in F$.
Then
\[
     \psi_0+t\phi
       =r_F+\sum_{l\in F}c_l
                    \left(\psi_l+\frac{t}{c_l}\phi_l\right)
 \tag{NP5}
\]
is an actual bounded continuous negative classical psh
function.  The normalization required by $\mathscr T_0$
also survives after decreasing $\epsilon_\phi$.
Indeed \textup{(UP5)} gives
\[
 T:=\sum_l c_lA_l=1-\sum_lc_l<1,
       \qquad\psi_0\ge T\rho\ge-T>-1.
\]
Choose in addition
$\epsilon_\phi\|\phi\|_\infty\le(1-T)/2$
(with no restriction if $\phi=0$).  Then
$-1<\psi_0+t\phi<0$.  It equals $\psi_0$ off a fixed compact set,
so its ordinary boundary value and compact strict negative
sublevels are unchanged.  Its top source equals that of
$\psi_0$ off the compact perturbation support by ordinary
bounded locality.  On a slightly larger compact set,
relative CLN gives finite mass.  Since $\psi_0$ has
finite total mass, so does the perturbed function.  All
finite positive power energies follow from boundedness.

For a fixed finite presentation, the Hessian bounds in
\textup{(NP4)}, the function-value bounds used for
strict negativity, and the finitely many $c_l$ are
uniform when $\|\phi\|_{2,\mathfrak a}\le1$.
The lower normalization uses the same fixed positive
margin $1-T$ and is uniform on that unit ball as well.
Scaling this fixed estimate proves \textup{(NP2)}.
The choices depend on the actual test presentation but
not on $S$ or its remaining degree.  This proves all
claims, including those for every further actual refinement.
\end{proof}

\begin{definition}[The normal probe defect]
Suppose $\ell\in\mathbb D_S^*$ and $u\in\mathscr C$
satisfy the variational normal relations
\[
 n=\ell-J(u),\qquad n(a)\le0\ (a\in\mathscr C),
                       \qquad n(u)=0.
 \tag{NP6}
\]
For example these hold for the uniquely constructed
$u=\Pi(\ell)$ whenever $\ell\in\mathbb D_{S,\mathrm s}^*$.
Define the finite scalar
\[
        \Xi_{\psi_0}(\ell,u)
                   =J(u)(\psi_0)-\ell(\psi_0)=-n(\psi_0)\ge0.
 \tag{NP7}
\]
The definition requires an existing normal pair.  It
does not assert existence of a variational maximizer for
every continuous functional on $\mathbb D_S$.
\end{definition}

\begin{theorem}[One scalar controls every original normal test]
\label{thm:r12-one-probe-normal-detector}
For every pair \textup{(NP6)} and every real original
compact smooth scalar $\phi$,
\[
            |n(\phi)|
                 \le\epsilon_\phi^{-1}\Xi_{\psi_0}(\ell,u).
 \tag{NP8}
\]
On a fixed finite presentation the right side may be
replaced by
$c_{\mathfrak a}^{-1}\max(1,\|\phi\|_{2,\mathfrak a})
\Xi_{\psi_0}(\ell,u)$; by homogeneity one obtains the
corresponding linear $C^2$-seminorm estimate for nonzero
tests.  This is a bound on original distributional test
pairings, not a total-variation norm assertion.

In particular $\Xi_{\psi_0}=0$ forces $n$ to vanish on
all original compact smooth scalar tests and on their
order-norm closed core $\mathbb L_{S,0}$ whenever $n$
extends continuously to that lattice.  It need not force
$n=0$ on $\mathbb D_S$ for a nonliteral energy dual.

Now let $\mu$ be a positive original Radon measure whose
\emph{literal} integration functional belongs to
$\mathbb D_{S,\mathrm s}^*$, and put $u=\Pi(\ell_\mu)$.
Then the following are equivalent:
\[
 \begin{split}
 \Xi_{\psi_0}(\mu)
   &:=\int_V(-\psi_0)\,d\widehat\mu_u
                      -\int_V(-\psi_0)\,d\mu=0;\\
 \widehat\mu_u&=\mu
        \quad\text{on all original compact smooth scalar tests};\\
 \text{there exists }v&\in\mathscr C
                         \text{ with }\widehat\mu_v=\mu.
 \end{split}
 \tag{NP9}
\]
If they hold, $v=u$ is unique.  For all these literal data,
\[
 \left|\int_V\phi\,d(\mu-\widehat\mu_u)\right|
                        \le\epsilon_\phi^{-1}\Xi_{\psi_0}(\mu).
 \tag{NP10}
\]
The same single geometry-only $\psi_0$ works for every
allowed arbitrary $S$, every remaining degree, every
datum in the printed class, and all actual refinements.
\end{theorem}

\noindent\textbf{Status: \texttt{NEW\_PROVED}} as the
single-probe reduction and its all-original quantitative
bound.  This is distinct from the Jensen-cluster probe
theorem: the defect here is a variational normal, not
a positive Jensen current.

\begin{proof}
Apply the polar inequality in \textup{(NP6)} to the
two actual cone members
$\psi_0\pm\epsilon_\phi\phi$ supplied by \textup{(NP1)}.
It gives
\[
               0\ge n(\psi_0)\pm\epsilon_\phi n(\phi),
\]
which is \textup{(NP8)}.  The presentation-wise bound
follows from \textup{(NP2)}.  If $n$ is continuous on
$\mathbb L_S$, vanishing on compact original smooth
tests extends to their defined closed core.  No density
of that core in the full lattice or difference space
is used.

For a literal datum, both integrals in \textup{(NP9)}
are finite by the energy-dual bound and mixed Holder.
The normal is literally
$n(\phi)=-\int\phi\,d(\mu-\widehat\mu_u)$,
so \textup{(NP8)} proves \textup{(NP10)}.
If the scalar defect is zero, every original smooth
test integral of the two Radon measures agrees.
Original scalar density and Radon uniqueness give
equality as original measures.  Conversely equality
of these measures gives equality of their finite
integrals against the actual continuous weight
$-\psi_0$, so the scalar defect is zero.

Any source potential $v$ has $J(v)=\ell_\mu$ by
the literal energy extension.  It is the unique
variational maximizer by the frozen convex energy
support inequality and strict separation.  Thus it
equals $u$.  This proves the third equivalence and
uniqueness.  Neither equality on additional GK test
forms nor an actual ambient classical representative
of every completed potential is inferred.
\end{proof}

\begin{corollary}[An additive strictness criterion for the source equation]
\label{cor:r12-additive-probe-source-criterion}
Under the literal datum hypotheses of \textup{(NP9)},
suppose its variational maximizer admits a decomposition
\[
                  u=v+\varepsilon\psi_0,
                  \qquad v\in\mathscr C,\quad\varepsilon>0.
 \tag{NP11}
\]
Then $\widehat\mu_u=\mu$ on all original tests.

The hypothesis is an \emph{additive cone decomposition}.
The pointwise inequalities $u\ge\varepsilon\psi_0$ or
$u\le\varepsilon\psi_0$ do not imply it.  Nor is it
enough that $u$ merely has a bounded classical lower
barrier.  Those replacements are contradicted by the
explicit incompatible bounded-source examples.
\end{corollary}

\noindent\textbf{Status: \texttt{NEW\_PROVED}} conditional
on the explicit geometric variation hypothesis, not on
an assertion restating the source equation.

\begin{proof}
The normal relation gives
$0=n(u)=n(v)+\varepsilon n(\psi_0)$.
Both terms are nonpositive, so $n(\psi_0)=0$.
The one-probe theorem applies.  Equivalently,
\textup{(NP1)} shows directly that every original
compact smooth direction can be varied with both
signs at $u$ within the cone; the normal must
annihilate it.  This is the step unavailable at
an arbitrary maximizer.
\end{proof}

\begin{proposition}[Density of the test-interior source family]
\label{prop:r12-dense-test-interior-sources}
For every $a\in\mathscr C$ and $\varepsilon>0$,
$a+\varepsilon\psi_0$ admits both signs of every
original compact smooth direction within $\mathscr C$.
These potentials approach $a$ in the energy-difference
norm as $\varepsilon\downarrow0$.  Their curvature
functionals satisfy the explicit estimate
\[
 \|J(a+\varepsilon\psi_0)-J(a)\|_*
       \le (r(a)+\varepsilon r(\psi_0))^q-r(a)^q.
 \tag{NP12}
\]
Thus this test-interior source family is dense in
$J(\mathscr C)$ for the signed operator dual norm.
It is not asserted dense among all positive data,
nor is an arbitrary prescribed datum automatically
the source of a test-interior potential.
\end{proposition}

\noindent\textbf{Status: \texttt{DIRECT\_COROLLARY}}
of the two-sided probe and the bounded mixed energy
form on $\mathbb D_S$.

\begin{proof}
For a fixed test $\phi$, use
$a+\varepsilon(\psi_0+t\phi)\in\mathscr C$
with $|t|\le\epsilon_\phi$.  The difference from
$a$ has norm $\varepsilon r(\psi_0)$.
Expand the $q$ slots of $J(a+\varepsilon\psi_0)$.
Each mixed functional has operator norm at most
the product of the roots of its differentiated
slots, by the proved mixed energy dual bound.
The binomial sum is precisely the right side of
\textup{(NP12)}.  This also proves the density
claim without approximating the original background.
\end{proof}

\begin{proposition}[Stability and a scalar lower bound for image distance]
\label{prop:r12-normal-probe-stability}
For $\ell,m\in\mathbb D_{S,\mathrm s}^*$ let
$\Xi(\ell)=\Xi_{\psi_0}(\ell,\Pi(\ell))$,
$\delta=\|\ell-m\|_*$ and
$R=\max(\|\ell\|_*,\|m\|_*)$.  Then
\[
            |\Xi(\ell)-\Xi(m)|
            \le r(\psi_0)\bigl(\delta+\sqrt{2qR\delta}\bigr).
 \tag{NP13}
\]
If $r(\psi_0)>0$, the normal satisfies
$\|\ell-J(\Pi(\ell))\|_*\ge\Xi(\ell)/r(\psi_0)$.
For $\|\ell\|_*>0$ this gives
\[
 \operatorname{dist}_*(\ell,J(\mathscr C))
 \ge
 \left(\sqrt{q\|\ell\|_*+\Xi(\ell)/r(\psi_0)}
                        -\sqrt{q\|\ell\|_*}\right)^2.
 \tag{NP14}
\]
Thus a strictly positive scalar normal defect certifies
a quantitative failure of source-image approximation.
The converse zero-defect assertion uses literal data
as in \textup{(NP9)}, not an arbitrary dual shadow.
\end{proposition}

\noindent\textbf{Status: \texttt{DIRECT\_COROLLARY}}
of the scalar definition and the proved nonlinear
retraction estimate \textup{(NI11)--(NI12)}.

\begin{proof}
Apply the dual norm to the difference of the two
values at $\psi_0$, then use \textup{(NI11)}.
This proves \textup{(NP13)}.  Testing the residual
norm at $\psi_0$ gives its displayed lower bound.
The right side of \textup{(NI12)}'s lower estimate
is increasing in that residual norm, yielding
\textup{(NP14)}.  No absolute-variation norm
$\delta_1$ is substituted for the signed dual norm.
\end{proof}

\paragraph{Why literal data cannot be omitted from the equivalence.}
In the radial example of Section~\ref{sec:r12-radial-pole-order-dual},
the functional $L_\infty$ vanishes on every bounded
lattice element and is nonnegative on $\mathscr C$.
Use the geometric bounded probe for that same background
and set $\ell=J(\psi_0)-L_\infty$, $u=\psi_0$.
Then $n=-L_\infty$ is nonpositive on the potential
cone and $n(u)=0$.  The convex support inequality
shows that $u$ is its unique variational maximizer.
The scalar defect is zero, but $n\ne0$ on the
energy-difference space.  Its restriction to every
original compact smooth test is zero, as the theorem
predicts.  This is a \texttt{NEW\_PROVED} example
of the nonliteral residual left undetected by a
current-test criterion.  It does not contradict
\textup{(NP9)} and does not assert strong continuity
of this $\ell$.

\paragraph{Remaining minimal issue.}
For literal finite-energy data the full original source
equation is now reduced to a single nonnegative
geometry-weighted variational defect.  Its vanishing
is still a genuine source-compatibility restriction:
the positive-capacity atom in \textup{(RF18)} has
strictly positive defect by \textup{(NP9)}.  A useful
general source theorem must prove vanishing from
additional prescribed-data or geometric hypotheses.
The probe theorem neither claims such unrestricted
vanishing nor changes OPEN-GQ, OPEN-CTO, or the
ordinary/extra-delta-test distinction.
